\paragraph{One signed coefficient calculation.}
\label{j:local:atlas}\label{j:local:moments}
All bounds in this paragraph concern one actual marginal. Put
$a_\pm=(1\pm\beta)/2$, $b_\pm=(\rho\pm\gamma)/2$ and
$m=\E X_+=\E X_->0$. The moment matrices below give
$m^2\le a_+a_-$ and $b_\pm m\ge a_\pm^2$, hence
\[
 \frac{(1+\beta)^{3/2}}{\sqrt{1-\beta}}-\rho\le\gamma
 \le\rho-\frac{(1-\beta)^{3/2}}{\sqrt{1+\beta}},\qquad
 m\ge\max\left\{\frac{1+\beta^2}{2\rho},\frac{a_+^2}{b_+},\frac{a_-^2}{b_-}\right\}.
\]
The Gram matrix of $X,|X|,X|X|$ similarly gives
$\mu_4\ge\rho^2+(\gamma-\rho\beta)^2/(1-\beta^2)$.
Integrating $(x-a)^2(R-x)\ge0$ also gives
$2m\le(R-\rho+2a+Ra^2)/[a(2R+a)]$ for $a>0$; the cover uses
$a=2/10,\ldots,12/10$.

The same matrices give both signed constant deductions at once.
Write $\delta=\gamma-\rho\beta$ and
\[
 P_s=\frac{(\rho+\gamma)(\delta_+)^2}{8(1+\beta)^2}
       +\frac{(\rho-\gamma)((-\delta)_+)^2}{8(1-\beta)^2},\qquad
 P_g=\frac{\delta^4}{4\rho(1-\beta^2)^2}.
\]
Then $\rho\mu_4/2-\eta_5/4\le\rho^3/4-\max(P_s,P_g)$.
Indeed, on either sign, $c=\E X_s^4\ge b^2/a$ and
$f=\E X_s^5\ge c^2/b$, so
$\rho c/2-f/4\le\rho^2b/4-(c-\rho b)^2/(4b)$.
Distance to $c\ge b^2/a$ gives $P_s$; using
$\eta_5\ge\mu_4^2/\rho$ and the Gram bound gives $P_g$.
The maximum is a choice of two deductions, not their sum.
These moment statements allow an atom at zero. The following central-gap
clauses apply only after $(-A,B)$, $A,B>0$, has been acquired as below.
